\section{Task and Evaluation Contract}
\label{sec:task_contract}

\subsection{Inputs and decision operations}
\label{sec:task_contract-inputs-and-decision-operations}

An inference input is \(x=(q,o,e,c,a_0)\), where \(q\) contains the current question and
patient information, \(o\) contains native options when present, \(e\) is the required
supplied evidence, \(c\) is the original system's allowed context, and \(a_0\) is its
saved answer when available. An explicit executable model is an additional optional input.
Each arm receives the same legal input for a given model configuration. Gold answers,
evaluation labels, classification explanations, hidden paired endpoints, and source
identifiers that could provide an answer shortcut are excluded from inference. Provenance
and scoring information remain in the offline evaluation ledger.

The benchmark uses five overlapping labels attached to specific judgments. R1 denotes
newly established support; R2 denotes withdrawn or contradicted support; R3 denotes a need
to compare competing judgments again; R4 denotes a consequence that must be inferred from
changed premises; and R5 denotes a judgment that should remain unaffected. Support changes
do not imply that the final answer must flip. R1--R3 may coexist, and R5 is established by
the task's meaning rather than by model agreement. The labels describe evaluation strata;
routing never reads them.

The system operates on each visible endpoint separately. It is not given both members of a
hidden evaluation pair and asked to recover the expected relation. Pair-level measurements
are constructed only after native predictions have been scored.

The project protocol permits external MedCorp retrieval while retaining all required
source evidence. In particular, MedCounterFact inputs preserve their stipulated evidence,
and MedRGB inputs retain their six fixed evidence-ratio conditions while also permitting
external retrieval. The latter is a project adaptation of MedRGB, rather than a
reproduction of its original static-evidence-only setting. Original and modified endpoints
have the same information-access permissions. The present study reuses the saved contexts
under that protocol.

\subsection{Native units and denominators}
\label{sec:task_contract-native-units-and-denominators}

The frozen v5 evaluation contains 13,905 unique inference inputs per model configuration,
15,616 native scoring mappings, and 6,104 non-reference evaluation units from 4,766
families. One input can support multiple mappings, and one unit can contain several native
tasks or variants. Reference inputs support paired evaluation but are excluded from the
main ALL score. These counts are not interchangeable sample sizes, and none is described
as a count of independently recruited patients.

Construction retained every eligible R1--R4 target in the preceding classified collection.
For R5, 450 eligible comparison units were randomly sampled without replacement from each
of eight source strata, using a seed generated once before selection and sorted unit
identifiers. Model predictions and prior v3 membership did not determine these draws. The
resulting R5 evaluation scope has 3,600 units. Frozen evaluation membership is distinct
from the original semantic labels: an unsampled R5 label attached to a retained R1 unit
does not automatically add that unit to the R5 score.

Let \(J_u\) be the non-reference native scoring records belonging to unit \(u\), and let
\(z_{muj}\in\{0,1\}\) be native main-answer correctness for arm \(m\). We use the frozen
unit score and overall mean

\begin{equation}
\label{eq:unit_score}
s_{mu}=\frac{1}{|J_u|}\sum_{j\in J_u}z_{muj},
\qquad
S_m=\frac{1}{6104}\sum_{u=1}^{6104}s_{mu}.
\end{equation}

The main metric is \(100S_m\). It is not a source macro-average, an average of the five
R-category scores, or input-level accuracy. Native rules include exact option keys, exact
sets for the primary multi-select score, normalized relation labels, and a frozen
canonical diagnostic mapping. Multi-select overlap and robustness auxiliary fields are
separate diagnostics. The diagnostic mapping is a local scoring contract, not the original
source's semantic language-model grader. For R4 clinical scenarios, correctness means
agreement with the source's expected answer.

The R-category denominators are R1: 2,866; R2: 79; R3: 733; R4: 13; and R5: 3,600. Because
these groups overlap, their counts and scores cannot be summed to recover ALL. The five
input formats contain 3,560 diagnosis inputs, 167 multi-select inputs, 989 relation
inputs, 2,700 robustness inputs with auxiliary output, and 6,489 other single-choice
inputs.

The earlier terminal-head experiments report the complete R1, R2, and R3 strata. Each
unit in these three strata has one non-reference scoring task, so their reported correct
counts divided by 2,866, 79, and 733 equal their corresponding category means. All 733 R3
units belong to R1; R1 and R2 intersect in four units, including three R3 units. Thus,
3,678 category memberships represent 2,941 unique units. Seven system evaluations of a
shared unit are repeated system observations, not seven independent patients. These
category-specific results do not define a seven-system CF-MoA ALL score.

\subsection{Paired correctness and uncertainty}
\label{sec:task_contract-paired-correctness-and-uncertainty}

The frozen paired ledger contains 6,838 maintain pairs and 1,780 respond pairs. Their
expected relations are based on the native task and reference answers, with compatible
answer meanings, rather than inferred mechanically from R labels. For pair \(p\) with
endpoints \(x_p^0\) and \(x_p^1\), the main paired score is

\begin{equation}
\label{eq:paired_score}
t_{mp}=\mathbb{1}[a_m(x_p^0)=y_p^0]\,
       \mathbb{1}[a_m(x_p^1)=y_p^1].
\end{equation}

Thus, a model receives credit only when both endpoints are correct; repeating an incorrect
answer does not count as successful maintenance. Native scoring, rather than literal
string equality where inappropriate, determines each indicator.

For each later full-v5 CF-MoA contrast, 5,000 paired bootstrap samples draw families with replacement using
seed 240924. All compared arms share a resampled family sequence. Each replicate averages
the atoms contributed by the sampled families, preserving unit weights within families
instead of assigning equal weight to family means. ALL uses 6,104 units in 4,766 families;
maintain uses 6,838 pairs in 2,806 families; respond uses 1,780 pairs in 1,780 families.
We report 95\% percentile intervals. These intervals describe variation in the exposed
benchmark and do not correct for method selection, multiple comparisons, or previous data
access.

The seven-system baseline and terminal-head tables report their complete observed point
estimates. We do not attach the later CF-MoA bootstrap intervals to those earlier
comparisons or treat overlapping category improvements as independent replications.

Families group the related variants of an original case; cross-source groups with exactly
matching original questions and options are connected into a dependency group. This
procedure does not establish that all semantic near-duplicates have been identified.
